\documentclass[UTF8,a4paper,12pt]{ctexart}
\usepackage[margin=2.2cm]{geometry}
\usepackage{graphicx}
\usepackage{hyperref,booktabs,longtable,tabularx,amsmath,amssymb}
\hypersetup{colorlinks=true,urlcolor=blue,linkcolor=black}
\title{视频生成模型发展讲义\\从扩散模型到可交互世界模型}
\author{整理日期：2026 年 9 月 12 日}
\date{}
\begin{document}
\maketitle
\tableofcontents
\section{学习路线}
本讲义按“视频表示—扩散与 DiT—效率—多模态—世界模型”递进。先理解 3D-VAE 如何把像素压缩成时空 latent，再比较 DDPM 与 Flow Matching，最后说明 Genie 和 LingBot 为什么需要动作条件与物理奖励。

\section{视频表示与扩散基础}
早期 VideoGPT、DVD-GAN 在像素或离散 VQ-VAE token 上自回归，因果性好但序列很长。Make-A-Video、Imagen Video 将图像扩散迁移到时空 U-Net。Sora 之后，主流改为 3D-VAE 潜空间与 Diffusion Transformer（DiT）。设视频 latent 为 $z_0\in\mathbb R^{T\times H\times W\times C}$，DDPM 损失为
\[
\mathcal L_{\rm DDPM}=\mathbb E_{z_0,\epsilon,t}\|\epsilon-\epsilon_\theta(z_t,t,c)\|_2^2.
\]
Flow Matching 直接学习速度场：
\[
\mathcal L_{\rm FM}=\mathbb E_{t,z_0,z_1}\|v_\theta(z_t,t,c)-(z_0-z_1)\|_2^2.
\]
视频 token 数 $N=T H W$，全注意力复杂度 $O(N^2)$，是长视频的主要瓶颈。

\section{关键模型}
\subsection{Genie：从生成视频到生成环境}
Genie（DeepMind，2024）论文：\href{https://arxiv.org/abs/2402.15391}{arXiv:2402.15391}。约 11B 参数，包含时空 tokenizer、从相邻帧推断离散 latent action 的模型，以及根据历史状态和动作预测下一 latent 的自回归 dynamics。它逐帧 rollout 并接受键盘/鼠标控制，学习目标是环境转移而非一次性 T2V 画面。

\subsection{HunyuanVideo}
腾讯 HunyuanVideo（2024）论文：\href{https://arxiv.org/abs/2412.03603}{arXiv:2412.03603}，\href{https://github.com/Tencent-Hunyuan/HunyuanVideo}{代码}。13B DiT 使用因果 3D-VAE（时间、空间、通道压缩约 4、8、16 倍），MLLM 编码文本，Transformer 采用 dual-stream 到 single-stream：先分别调制视频和文本 token，后段拼接并做全注意力融合。联合图像/视频训练，支持 T2V、I2V、V2V 编辑和个性化。

\subsection{Wan 2.1}
Wan（2025）论文：\href{https://arxiv.org/abs/2503.20314}{arXiv:2503.20314}，\href{https://github.com/Wan-Video/Wan2.1}{代码}。提供 1.3B/14B，使用因果 Wan-VAE、DiT、Flow Matching、文本 cross-attention 和全时空注意力，覆盖 T2V、I2V、V2V 编辑、个性化等任务。1.3B 约 8.2GB 显存，14B 质量高但推理慢。

\subsection{MiniMax Hailuo-02 与 H3}
Hailuo-02 使用 Noise-aware Compute Redistribution（NCR）按噪声阶段分配计算，官方宣称约 2.5 倍效率提升（\href{https://www.minimax.io/news/minimax-hailuo-02}{官方介绍}）。MiniMax H3（2026-07-31，\href{https://www.minimax.io/blog/minimax-h3}{官方博客}）以 Contextual Omni Representation、H3-VAE、H3-Omni Transformer 和 In-Context Regeneration 统一文字、图像、视频、音频，最高 2K/15 秒并原生生成立体声音频。权重开放状态应以最新公告为准。

\subsection{SANA-Video}
SANA-Video（2025）论文：\href{https://arxiv.org/abs/2509.24695}{arXiv:2509.24695}。Linear DiT 和块式自回归使用恒定内存 KV 状态，支持分钟级 720$\times$1280；实测延迟约为 Wan-1.3B 的 $1/16$。SANA-Video 2.0 5B 采用 75\% 线性注意力、25\% softmax anchor、Attention Residual 和 LTX-2.3 VAE，并提供 DMD 四步蒸馏模型（\href{https://nvlabs.github.io/Sana/docs/sana_video2/}{文档}）。

\subsection{LingBot 与 FLUX.3}
LingBot-Video（2026）论文：\href{https://arxiv.org/abs/2607.07675}{arXiv:2607.07675}，采用 MoE DiT（仓库型号 30B-A3B）、机器人操作/导航数据和物理、任务完成度等多维奖励；LingBot-World 2.0 进一步支持动作交互与实时蒸馏（\href{https://arxiv.org/abs/2607.07534}{arXiv:2607.07534}）。FLUX.3 Video（2026-07，Black Forest Labs）目前主要有产品/API 资料，宣称约 20 秒、关键帧和原生音频；训练目标、VAE 与 Transformer 尚未公开，不能进行参数级复现比较。

\section{时间线与对比}
\begin{longtable}{p{2.2cm}p{4.2cm}p{7.2cm}}
\toprule 时间 & 模型 & 方法演进\\\midrule
2019--21 & VideoGPT、DVD-GAN & 离散 token/像素自回归与 GAN\\
2022 & Make-A-Video、Imagen Video & 时空 U-Net 扩散\\
2023 & VideoCrafter、AnimateDiff、Lumiere & 时空注意力、运动模块、DiT 基础\\
2024 & Sora、Genie、HunyuanVideo & 潜空间 DiT；latent-action 世界模型；13B 开源 DiT\\
2025 & Wan 2.1、Hailuo-02、SANA-Video & Flow Matching、多任务开源；NCR；Linear DiT\\
2026 & LingBot、SANA-Video2、MiniMax H3、FLUX.3 & MoE 具身、混合注意力与蒸馏、统一音视频与关键帧\\\bottomrule
\end{longtable}

\begin{longtable}{p{2.8cm}p{4.4cm}p{4.2cm}p{2.4cm}}
\toprule 模型 & 核心范式 & 条件/输出 & 开放性\\\midrule
Genie & latent action+AR dynamics & 动作控制、逐帧 rollout & 研究原型\\
HunyuanVideo & 13B DiT、双流到单流、3D-VAE & T2V/I2V/V2V，约 5 秒 & 开源\\
Wan 2.1 & DiT+Flow Matching、Wan-VAE & T2V/I2V/编辑 & 开源\\
Hailuo-02 & NCR & 6/10 秒、最高 1080p & 闭源\\
SANA-Video2 & 线性+softmax anchor、DMD & 5B、720p/8 秒 & 开源 5B\\
MiniMax H3 & Omni 表示、H3-VAE、原生音频 & 2K/15 秒、多参考 & 计划开放\\
LingBot-Video & MoE DiT、具身奖励 & 动作条件、T2V/TI2V & 开源\\
FLUX.3 & 统一图像/视频/音频/动作 & 20 秒、关键帧、音频 & 闭源\\\bottomrule
\end{longtable}

\section{榜单阅读原则}
截至 2026-09-12 没有跨任务公认的单一总榜。应分别报告闭源综合质量、开源质量、效率和交互世界模型，并记录评测日期、分辨率、时长、是否含音频及人工/自动指标。建议将 FLUX.3、Veo、Sora、MiniMax H3 视为闭源第一梯队候选，将 LingBot/Wan/Hunyuan 作为开源候选，SANA-Video2 作为效率候选，LingBot-World 与 Genie 作为交互候选；名次需随公开榜单更新。



\section{问题驱动的发展脉络与论文原图}

本章按“瓶颈 $→$ 方法 $→$ 代价 $→$ 后续影响”阅读。插图均来自论文 arXiv 源码或作者官方仓库，未自行重绘。

\subsection{Genie：为什么要从生成视频转向生成环境？}
普通 T2V 模型优化的是整段视频的去噪/似然，用户只能在视频级提供文本；它没有显式的动作变量，因此难以回答“按下左键后下一帧会怎样”。Genie 的解决方案是三模块：时空视频 tokenizer、无监督 latent-action model（LAM）和自回归 dynamics。LAM 从相邻帧反推出离散动作，dynamics 以历史 token 与动作逐帧预测下一帧。代价是自回归 rollout 比一次性扩散慢、分辨率较低；影响是把视频模型引向世界模型和具身智能。

\begin{figure}[htbp]
\centering

\includegraphics[width=.94\linewidth]{figures/original/genie/genie_architecture.png}
\caption{Genie 原论文架构图：tokenizer、latent-action model 与 dynamics model。来源：\href{https://arxiv.org/html/2402.15391}{arXiv HTML，Fig.2/3/6}。}
\end{figure}

\subsection{HunyuanVideo：如何让开源 Video DiT 接近闭源质量？}
瓶颈从“能生成”变为高分辨率下的 token 数、文本对齐和训练资源。HunyuanVideo 采用 causal 3D-VAE 压缩时空 latent，使用 MLLM 编码文本，并在 DiT 中先分离视频/文本流、再合并为 single-stream 全注意力。双流阶段保留模态专门化，单流阶段增强深层融合；代价是 13B 模型显存和训练成本高。

\begin{figure}[htbp]
\centering
\includegraphics[width=.96\linewidth]{figures/original/hunyuanvideo_overall.png}
\caption{HunyuanVideo 官方架构图（Causal 3D-VAE、文本编码器与 DiT）。来源：\href{https://github.com/Tencent/HunyuanVideo/blob/main/assets/hunyuanvideo.pdf}{Tencent HunyuanVideo 官方仓库}。}
\end{figure}

\subsection{Wan 2.1：如何把 T2V、I2V 与编辑统一到一个开放基线？}
HunyuanVideo 的高质量伴随较高部署门槛，研究者还需要一个多任务、可消费级运行的开放模型。Wan 以 Wan-VAE、DiT 和 Flow Matching 统一 T2V/I2V/V2V；文本或图像条件通过 cross-attention 注入，Flow Matching 用 ODE 速度场替代 DDPM 多步噪声预测。1.3B 版本降低显存门槛，14B 版本追求质量；但全时空注意力仍让长视频昂贵。

\begin{figure}[htbp]
\centering
\includegraphics[width=.94\linewidth]{figures/original/wan/video_dit/dit_arch.pdf}
\caption{Wan 原论文 Video DiT 架构图。来源：\href{https://arxiv.org/html/2503.20314}{arXiv HTML，Fig.9/10}。}
\end{figure}

\subsection{SANA-Video：当 token 数成为显存瓶颈，如何生成分钟级视频？}
标准 self-attention 的 $O(N^2)$ 成本随帧数迅速增长。SANA-Video 将注意力重写为线性形式，并用 block-wise causal KV cache 保存固定大小的状态；SANA-Video 2.0 再混入 softmax anchor 恢复部分精细对应，并用 DMD 蒸馏到 4 步。它牺牲部分全局精确匹配，换取长视频和低延迟。

\begin{figure}[htbp]
\centering
\includegraphics[width=.94\linewidth]{figures/original/sana/model.pdf}
\caption{SANA-Video 原论文模型总览。来源：\href{https://arxiv.org/html/2509.24695}{arXiv HTML，Fig.2}。}
\end{figure}

\subsection{LingBot-Video：如何把“好看”推进到“可执行”？}
面向机器人时，仅有互联网视频的审美和文本对齐不够，还需要物理合理性、动作语义和任务完成度。LingBot-Video 使用稀疏 MoE DiT，让每个 token 只激活少量专家；数据混合互联网视频、操作/导航/第一视角轨迹；后训练加入物理与任务奖励。它保留扩散模型的画质，同时比 dense 大模型更省激活计算，但训练、路由稳定性和评测更加复杂。

\begin{figure}[htbp]
\centering
\includegraphics[width=.94\linewidth]{figures/original/lingbot/architecture/architecture.pdf}
\caption{LingBot-Video 原论文架构图。来源：\href{https://arxiv.org/html/2607.07675}{arXiv HTML，Fig.2}。}
\end{figure}

\subsection{如何比较这些图？}
沿三条轴阅读：表示从离散 token 到连续 VAE latent；动力学从 AR 下一帧到扩散/Flow ODE；计算组织从 dense attention 到 linear attention 和 MoE。Genie 最强调动作因果性，HunyuanVideo/Wan 最强调通用视觉质量，SANA 最强调内存和速度，LingBot 最强调物理与任务可执行性。

\section{原图资产索引}
\begin{itemize}
\item Genie：\texttt{figures/original/genie/genie\_architecture.png}。
\item HunyuanVideo：\texttt{figures/original/hunyuanvideo\_overall.png}。
\item Wan：\texttt{figures/original/wan/video\_dit/dit\_arch.pdf}。
\item SANA-Video：\texttt{figures/original/sana/model.pdf}。
\item LingBot-Video：\texttt{figures/original/lingbot/architecture/architecture.pdf}。
\end{itemize}
\section{世界模型延伸阅读}
\begin{itemize}
\item \href{https://www.xunhuang.me/blogs/world_model.html}{Xun Huang: From Video Generation to Video World Models}：区分视觉逼真、交互、规划和记忆，并讨论视频生成与 3D 世界模型的互补关系。
\item \href{https://video-world-models.github.io/blog/video-world-models/}{Video World Models}：解释双向扩散的因果局限，以及 Self-Forcing、检索增强和长上下文记忆。
\item \href{https://arxiv.org/abs/2505.20171}{Long-Context State-Space Video World Models}：用 block-wise SSM 扩展时间记忆。
\item \href{https://arxiv.org/abs/2505.21996}{Learning World Models for Interactive Video Generation}：分析误差累积与全局状态条件。
\item \href{https://arxiv.org/abs/2608.30692}{Can Video World Models Track Unobserved World States?}：从不可见状态外推检验世界模型是否真正有记忆。
\end{itemize}
\
\section{2025--2026 扩展：代表性 SOTA 与架构进展}

\subsection{Step-Video-T2V：深压缩 VAE 与视频 DPO}
\begin{figure}[htbp]\centering\includegraphics[width=.94\linewidth]{figures/original/recent/step_video_model_architecture.png}\caption{Step-Video-T2V 原论文模型架构。来源：\href{https://arxiv.org/abs/2502.10248}{arXiv:2502.10248}。}\end{figure}
30B 级 Video DiT 的瓶颈是时空 token 数和伪影。Step-Video 采用 8 倍时间、16x16 空间压缩的 Video-VAE、双语文本编码器、3D full-attention DiT、Flow Matching 与 Video-DPO；代价是复杂物理交互仍有限。

\subsection{HunyuanVideo 1.5：选择性滑窗注意力与级联超分}
\begin{figure}[htbp]\centering\includegraphics[width=.94\linewidth]{figures/original/recent/hunyuanvideo15_architecture.png}\caption{HunyuanVideo 1.5 原论文架构。来源：\href{https://arxiv.org/abs/2511.18870}{arXiv:2511.18870}。}\end{figure}
HunyuanVideo 1.5 用 selective/sliding tile attention 减少无效交互，基础 DiT 后接视频超分网络；它回答了高清质量与推理速度的矛盾。

\subsection{Epona：视频世界模型与轨迹规划解耦}
\begin{figure}[htbp]\centering\includegraphics[width=.94\linewidth]{figures/original/recent/epona_pipeline.png}\caption{Epona 原论文流程图。来源：\href{https://arxiv.org/abs/2506.24113}{arXiv:2506.24113}。}\end{figure}
Epona 用 causal transformer 建模 latent dynamics，再由两个扩散分支分别生成未来视频和车辆轨迹，并用 chain-of-forward 训练缓解 rollout 误差；它把渲染与规划放入同一世界模型。

\subsection{SkyReels V2--V4：长片、参考一致性和音视频统一}
SkyReels-V2 引入 shot-aware caption、motion RL 与 diffusion forcing；V3 统一参考图 I2V、延长和音频引导；V4 使用 dual-stream MMDiT，将生成、修补和编辑统一为条件 inpainting。论文：\href{https://arxiv.org/abs/2504.13074}{V2}、\href{https://arxiv.org/abs/2601.17323}{V3}、\href{https://arxiv.org/abs/2602.21818}{V4}。

\subsection{LanDiff、FlashVideo 与 FSVideo：语义先行和极限加速}
LanDiff 用高压缩 1D 语义 token 让语言模型先规划、扩散模型后精修；FlashVideo 用低分辨率大模型保证语义与运动，再用约 4 步 Flow Matching 补高清细节；FSVideo 在高度压缩 latent 中加入 layer memory 和多分辨率 refiner。论文：\href{https://arxiv.org/abs/2503.04606}{LanDiff}、\href{https://arxiv.org/abs/2502.05179}{FlashVideo}、\href{https://arxiv.org/abs/2602.02092}{FSVideo}。

\subsection{长视频稳定性：Endless World、Cycle-World、BAgger 与 Diagonal Distillation}
这些工作都处理自回归 exposure bias 与误差累积：Endless World 注入全局 3D 特征，Cycle-World 用反向预测和 cycle guidance，BAgger 从自身错误 rollout 聚合纠正轨迹，Diagonal Distillation 让早期 chunk 使用更多去噪步、后续 chunk 使用更少步。论文：\href{https://arxiv.org/abs/2512.12430}{Endless World}、\href{https://arxiv.org/abs/2607.11836}{Cycle-World}、\href{https://arxiv.org/abs/2512.12080}{BAgger}、\href{https://arxiv.org/abs/2603.09488}{Diagonal Distillation}。
\section{Xun Huang 博客的统一视角}
Xun Huang 将视频世界模型按能力分为 faithfulness、interactivity、planning 与 statefulness 四层：画面逼真只是起点，模型还要能接受动作、评估轨迹并维护跨长时/多视角的隐藏状态。因此，VBench 或人类美学分数不能单独证明世界模型能力。

架构演进也可归纳为三条线：causal/self-forcing 将双向教师转为可流式 rollout；SSM、线性注意力、KV cache 与显式 world-state 解决长上下文；3D 几何、相机和动作条件把像素预测推进到可执行模拟。Genie、SANA-WM、LingBot、Epona、Endless World 和 WorldWeaver 分别代表这些取舍。博客：\href{https://www.xunhuang.me/blogs/world_model.html}{From Video Generation to Video World Models}。
\section*{文献检索说明}
2025--2026 扩展部分使用 DeepXiv CLI 的 \texttt{deepxiv search} 语义检索，并对入选论文用 arXiv HTML/源码核对方法与架构图。
\end{document}




